\documentclass[11pt]{article}

\usepackage[margin=1in]{geometry}
\usepackage{amsmath,amssymb,booktabs,tabularx}

\title{Martingale tilt equation}
\author{}
\date{}

\begin{document}
\maketitle

\noindent\textbf{Status.} Draft\\
\textbf{Plan item.} Martingale tilt equation for \(\theta\)\\
\textbf{Depends on.} \texttt{08\_pricing\_density\_normalization.tex}\\
\textbf{Feeds.} Risk-neutral characteristic function, terminal pricing measure, martingale diagnostics

\section{Purpose}

Derive the scalar equation that determines the terminal return tilt
\(\theta\).
The activity premia \(\pi_+\), \(\pi_-\), and signed-imbalance exposure risk
price \(\rho_\Omega\) are pricing-kernel coordinates, while \(\theta\) is
solved so that the discounted ex-dividend terminal price is a martingale.

\section{Setup}

Write the terminal log return as
\[
    R_{t,T}
    =
    d_{t,T}
    +A_{t,T}^{\Omega}
    +X_{t,T}^{\mathrm{core}},
\]
where \(d_{t,T}\) is deterministic conditional on \(\mathcal F_t\) and
\(A_{t,T}^{\Omega}+X_{t,T}^{\mathrm{core}}\) is the stochastic Janiform
terminal component whose cumulant is \(K_{t,T}\).
\(A_{t,T}^{\Omega}\) is the unweighted market-pressure anchor-pressure
movement, while \(X_{t,T}^{\mathrm{core}}\) is the mean-reversion-weighted sum
of interval core shocks.
Let \(B_{t,T}^{\Omega}\) be the signed-balance exposure integral satisfying
\(A_{t,T}^{\Omega}=\gamma_\Omega B_{t,T}^{\Omega}\).
Under the normalized density from
\texttt{08\_pricing\_density\_normalization.tex},
\[
    \frac{dQ^{\theta,\pi,\rho_\Omega}}{dP}\bigg|_{\mathcal F_T}
    =
    M_T(\theta,\pi_+,\pi_-,\rho_\Omega).
\]

\section{Notation}

\begin{tabularx}{\textwidth}{lXX}
\toprule
Symbol & Meaning & Units or domain\\
\midrule
\(S_t\) & Valuation-time spot price & Strictly positive\\
\(S_T\) & Terminal price, \(S_T=S_t e^{R_{t,T}}\) & Strictly positive\\
\(r\) & Continuously compounded risk-free rate & Annualized real\\
\(q\) & Continuously compounded dividend yield & Annualized real\\
\(\tau\) & Year fraction from \(t\) to \(T\) & Positive\\
\(\theta\) & Scalar martingale tilt & Real transform domain\\
\(\pi_\pm\) & Activity-premium multipliers & Strictly positive\\
\(\rho_\Omega\) & Signed-imbalance exposure risk price & Real finite\\
\(B_{t,T}^{\Omega}\) & Signed-balance exposure integral & Real finite\\
\(K_{t,T}\) & Physical stochastic terminal cumulant & Finite at requested arguments\\
\bottomrule
\end{tabularx}

\section{Assumptions used}

\begin{tabularx}{\textwidth}{XXX}
\toprule
Assumption & Where it is used & Why it is needed\\
\midrule
The terminal price is \(S_T=S_t e^{R_{t,T}}\).
& Translating the martingale condition into a first exponential moment.
& Option pricing is performed on terminal prices, not only log returns.\\
The pricing density is normalized.
& Rewriting \(Q\)-expectations as ratios of \(P\)-moments.
& The denominator cumulant cancels total probability.\\
\(K_{t,T}\) is finite at both \((\theta,\rho_\Omega)\) and \((\theta+1,\rho_\Omega)\).
& Evaluating the density denominator and the first price moment numerator.
& The martingale equation requires both moments.\\
The deterministic terminal component is known at valuation time.
& Adding \(d_{t,T}\) outside the cumulant difference.
& Deterministic log returns shift the log martingale equation directly.\\
\bottomrule
\end{tabularx}

\section{Result}

\noindent\textbf{Identity.} The martingale tilt \(\theta\) is any admissible
real solution of
\[
    d_{t,T}
    +K_{t,T}(\theta+1,\log\pi_+,\log\pi_-,\rho_\Omega)
    -K_{t,T}(\theta,\log\pi_+,\log\pi_-,\rho_\Omega)
    =
    (r-q)\tau.
\]
The model uses this equation to derive \(\theta\).
It is not an independent free pricing parameter.

\section{Derivation}

The ex-dividend martingale condition is
\[
    \mathbb E^{Q^{\theta,\pi,\rho_\Omega}}\!\left[\frac{S_T}{S_t}\mid\mathcal F_t\right]
    =
    e^{(r-q)\tau}.
\]
Since
\[
    \frac{S_T}{S_t}
    =
    e^{R_{t,T}}
    =
    e^{d_{t,T}}
    e^{A_{t,T}^{\Omega}+X_{t,T}^{\mathrm{core}}},
\]
the terminal expectation factorizes as
\[
    \mathbb E^{Q^{\theta,\pi,\rho_\Omega}}\!\left[e^{R_{t,T}}\mid\mathcal F_t\right]
    =
    e^{d_{t,T}}
    \mathbb E^{Q^{\theta,\pi,\rho_\Omega}}\!\left[
        e^{A_{t,T}^{\Omega}+X_{t,T}^{\mathrm{core}}}
        \mid\mathcal F_t
    \right].
\]

Using the density \(M_T\), the stochastic expectation is
\[
    \mathbb E^{Q^{\theta,\pi,\rho_\Omega}}\!\left[
        e^{A_{t,T}^{\Omega}+X_{t,T}^{\mathrm{core}}}
        \mid\mathcal F_t
    \right]
    =
    \mathbb E^P\!\left[
        e^{A_{t,T}^{\Omega}+X_{t,T}^{\mathrm{core}}}
        M_T(\theta,\pi_+,\pi_-,\rho_\Omega)
        \mid\mathcal F_t
    \right].
\]
Substituting the normalized density gives
\[
    \mathbb E^{Q^{\theta,\pi,\rho_\Omega}}\!\left[
        e^{A_{t,T}^{\Omega}+X_{t,T}^{\mathrm{core}}}
        \mid\mathcal F_t
    \right]
    =
    \frac{
        \mathbb E^P\!\left[
            \exp\!\left(
                (\theta+1)A_{t,T}^{\Omega}
                +(\theta+1)X_{t,T}^{\mathrm{core}}
                +\log(\pi_+)N_{t,T}^+
                +\log(\pi_-)N_{t,T}^-
                -\rho_\Omega B_{t,T}^{\Omega}
            \right)
            \mid\mathcal F_t
        \right]
    }{
        \exp\!\left(K_{t,T}(\theta,\log\pi_+,\log\pi_-,\rho_\Omega)\right)
    }.
\]
By definition of \(K_{t,T}\), this is
\[
    \exp\!\left(
        K_{t,T}(\theta+1,\log\pi_+,\log\pi_-,\rho_\Omega)
        -K_{t,T}(\theta,\log\pi_+,\log\pi_-,\rho_\Omega)
    \right).
\]
Therefore
\[
    \mathbb E^{Q^{\theta,\pi,\rho_\Omega}}\!\left[e^{R_{t,T}}\mid\mathcal F_t\right]
    =
    \exp\!\left(
        d_{t,T}
        +K_{t,T}(\theta+1,\log\pi_+,\log\pi_-,\rho_\Omega)
        -K_{t,T}(\theta,\log\pi_+,\log\pi_-,\rho_\Omega)
    \right).
\]
Equating this to \(e^{(r-q)\tau}\) and taking logs gives
\[
    d_{t,T}
    +K_{t,T}(\theta+1,\log\pi_+,\log\pi_-,\rho_\Omega)
    -K_{t,T}(\theta,\log\pi_+,\log\pi_-,\rho_\Omega)
    =
    (r-q)\tau.
\]
This is the martingale tilt equation solved for \(\theta\).

\section{Domain and validity}

The equation is valid only for real \(\theta\) such that both
\(\theta\) and \(\theta+1\) are inside the terminal cumulant domain for the
given \(\pi_+\), \(\pi_-\), and \(\rho_\Omega\).
The density must be normalized, the activity premia must be strictly positive,
the pressure risk price must be finite,
and the affine schedule evaluations must remain finite.
Existence and uniqueness are numerical root-solving questions
inside the admissible interval and are not implied by normalization alone.

\section{Consequences}

\begin{itemize}
    \item The martingale tilt is derived from no-arbitrage and the chosen pricing-kernel coordinates.
    \item The same denominator cumulant used for density normalization appears in the martingale equation.
    \item The numerator uses a \(+1\) shift because the martingale condition is a first price moment.
    \item Martingale diagnostics should report the residual of this equation, not a fitted \(\theta\) objective.
\end{itemize}

\section{Checks}

\begin{tabularx}{\textwidth}{XXX}
\toprule
Check & Expected result & Evidence location\\
\midrule
Residual equation
& The residual is \(d_{t,T}+K(\theta+1)-K(\theta)-(r-q)\tau\).
& \texttt{evaluate\_terminal\_affine\_martingale\_residual} in \texttt{crates/janiform-engine-pricing/src/spectral/pricer/terminal\_affine/core.rs}\\
Theta solve
& The terminal affine pricer solves \(\theta\) through a bracketed local-tilt residual solve.
& \texttt{solve\_terminal\_affine\_martingale\_theta} in \texttt{crates/janiform-engine-pricing/src/spectral/pricer/terminal\_affine/core.rs}\\
Martingale diagnostics
& The reported absolute martingale residual is within the configured tolerance.
& Test \texttt{terminal\_affine\_measure\_change\_normalizes\_characteristic\_and\_martingale}\\
\bottomrule
\end{tabularx}

\end{document}
